\scanpage{59}
\diagram{
\[
\begin{array}{ccc}
\begin{gathered}\text{\unsure{forme extérieure}}\\G_1\ \text{\unsure{donnée}}\end{gathered}&
\left\{\begin{gathered}
\begin{tikzcd}[row sep=large]
{H^1(S,G_1)\longrightarrow H^1(S,\Gamma_1)\quad\text{injectif ?}\quad\text{Non}}\arrow[d,Rightarrow]\\
{\begin{gathered}\deleted{\unsure{H^1}}\\\underline{\operatorname{Aut}}(G)\longrightarrow\underline{\operatorname{Ext}}(G)\quad\text{admet tjrs}\\\text{une section}\end{gathered}}\arrow[d,Leftrightarrow]\\
{\operatorname{Aut}(G)\longrightarrow\operatorname{Ext}(G)\quad\text{tjrs surjectif}}
\end{tikzcd}
\end{gathered}\right\}&
\begin{gathered}\text{faux en général sur un corps}\\\text{vrai pour les formes}\\\text{intérieures}\quad\rho=0\\\text{ou sur les anneaux}\\\text{de Tsen}\end{gathered}
\end{array}
\]
}

\diagram{
\[
\begin{array}{cc}
\left.\begin{gathered}
\begin{tikzcd}[row sep=large]
{H^1(S,B_1)\longrightarrow H^1(S,G_1)\quad\text{injectif}}\arrow[d,Leftrightarrow]\\
{\begin{gathered}\text{Deux sous-groupes de Borel \inserted{d'un $=$ $G$}}\\\text{tjrs conjugués}\end{gathered}}\arrow[d,Leftrightarrow]\\
{\begin{gathered}\text{Pour \uuline{tt.}}\ G\ \text{et}\ B,\quad\deleted{\text{le noyau de}}\\H^1(S,B)\longrightarrow H^1(S,G)\quad\deleted{\text{est injectif}}\end{gathered}}\arrow[d,Leftrightarrow]\\
{\begin{gathered}\text{Pour tt}\ G,\ \text{et}\ B,\quad\deleted{\text{le noyau de}}\\H^1(S,B)\mathbin{-}H^1(S,G)\quad\text{est \unsure{bijectif|injectif}}\end{gathered}}
\end{tikzcd}
\end{gathered}\right\}&
\begin{gathered}-\\\text{vrai pour les formes}\\\text{intérieures}\ (\rho=0)\\\text{ou sur les anneaux}\\\text{de Tsen.}\end{gathered}
\end{array}
\]
}

\margin{\unsure{d'un cas} \unsure{la} $=$ \unsure{chose pour} tt \unsure{extension de} $S$ \inserted{(ou $S'$ \uuline{local})} \unsure{la même chose} est équivalent à dire que $G\longrightarrow G/B$ est fibré \uline{localement trivial}.}

\diagram{
\[
\begin{array}{cc}
\left.\begin{gathered}
\begin{tikzcd}[row sep=large]
{H^1(S,N_1)\mathbin{-}H^1(S,G_1)\quad\text{injectif}}\arrow[d,Rightarrow]\\
{\begin{gathered}\text{Deux sous groupes de Cartan du $=$}\ G\\\text{tjrs conjugués}\end{gathered}}\arrow[d,Leftrightarrow]\\
{\begin{gathered}\text{Pour tt.}\ G\ \text{et}\ T,\ \text{le noyau de}\\H^1(S,N(T))\mathbin{-}H^1(S,G)\\\text{trivial}\end{gathered}}\arrow[d,Leftrightarrow]\\
{\begin{gathered}\rule{4cm}{0.3pt}\\\hspace{1cm}\rule{2cm}{0.3pt}\quad\text{injectif}\end{gathered}}
\end{tikzcd}
\end{gathered}\right\}&
\begin{gathered}\deleted{\text{\unsure{Souvent}}}\ \text{faux déjà dans}\\\text{les corps \unsure{non}}\\\text{\unsure{p-adiques} : \unsure{leurs}}\\\text{\unsure{clôtures parfaites},}\\\text{$=$ pour les formes}\\\deleted{\unsure R}\ \text{intérieures, .....}\end{gathered}
\end{array}
\]
}
\margin{58}
